\newcommand{\tipoRequerimento}{Solicitação de Histórico de Consumo e Faturamento}

\section{FUNDAMENTAÇÃO LEGAL DA REQUISIÇÃO}

\begin{enumerate}[resume]
  \item Conforme estabelece o art. 399 da Resolução Normativa nº 1.000/2021, é um direito do consumidor poder apresentar, a qualquer momento, pedidos de informação, reclamações, sugestões ou denúncias à distribuidora, utilizando para isso os canais de atendimento por ela disponibilizados.
  \Citacao{Art. 399. O consumidor e demais usuários podem requerer informações, solicitar e cancelar serviços, encaminhar reclamações, elogios, sugestões, denúncias e solicitar o encerramento contratual nos canais de atendimento disponibilizados pela distribuidora. (Redação dada pela REN ANEEL 1.042, de 20.09.2022)

  § 1º O consumidor e demais usuários podem requerer informações, esclarecer dúvidas, encaminhar sugestões, reclamações, denúncias, elogios e sugestões à Ouvidoria da distribuidora, se houver, à agência estadual conveniada ou, na inexistência desta, à ANEEL, observado o art. 421 e o art. 423.

  § 2º O atendimento deve ser classificado e registrado de acordo com opção do consumidor e demais usuários quanto à natureza de sua demanda, entre informação, reclamação, serviço, denúncia, elogio ou sugestão, conforme os procedimentos e tipologia estabelecidos pela ANEEL.

  § 3º No caso de atendimento por canal telefônico de demanda para a qual exista previsão regulatória para o envio de documentos, a distribuidora deve registrar a demanda e informar ao consumidor e demais usuários a relação de documentos e os canais para envio. (Incluído pela REN ANEEL 1.042, de 20.09.2022)

  § 4º No caso do §3º, a contagem do prazo para tratamento da demanda deve ficar suspensa até o recebimento dos documentos, podendo ser indeferida pela distribuidora se o recebimento não ocorrer em até 5 dias úteis, exceto no caso de ressarcimento de danos, de que trata o art. 605. (Incluído pela REN ANEEL 1.042, de 20.09.2022)
  }
  \item Por fim, instruem o presente expediente os documentos abaixo relacionados, que se juntam como anexos e passam a integrá-lo para todos os fins de direito.
\end{enumerate}

\section{DOS PEDIDOS}

\begin{enumerate}[resume]
\item 
Perante o exposto, solicita-se a concessionária que forneça o \textbf{histórico de consumo e faturamento} da unidade consumidora agregadora dos pontos de iluminação pública (IP) faturados por estimativa, sob a modalidade tarifária B4a (acervo municipal), referente aos últimos \textbf{114 meses}, contendo, no mínimo, a discriminação dos valores faturados em kWh e em reais (R\$).
\end{enumerate}
